% Generated by task tex-libraries from dossiers/lib-others.md, section E (and F.1). The Layer 0
% sources are quoted with `git show b435631:<path>`; the other snippets are copied from the dossier.
\begingroup\sloppy
\section{Layer 0 of the proof system, audited}\label{sec:gen-l0}

Layer 0 is the proof system's interface to the libraries: the committed column, the challenge
samplers, the size of the challenge space, the oracle interface of the stack, and two counting
bounds. It lives in \file{LeanerVM/Protocol/Field.lean}, \file{Basic.lean},
\file{ToArkLib/Oracles.lean} and \file{ToVCVio/UniformSample.lean}, with tests in
\file{tests/LeanerVMTests/Protocol/Field.lean}. Everything below is at \code{b435631}; what the
upgrade to \code{144c5aa} changes is said where it matters. The audit's five results:
\begin{enumerate}
\item \textbf{The samplers are uniform, by proof} (\cref{sec:gen-l0-samplers}). The instance
  diamond on $\K$ at the old pins is definitionally one term; the $\K$ sampler has no consumer
  other than the $\E$ sampler.
\item \textbf{The column oracle answers the specification's evaluation}: query $r \in \E^{\mu}$,
  answer $\mle{q}(r) = \sum_i \mathrm{ofK}(q_i)\,\eq(r,i)$, in the specification's little-endian
  cube order (\cref{sec:gen-l0-oracle}).
\item \textbf{\lean{card_E} is correct} and robust to the choice of instance
  (\cref{sec:gen-l0-card}).
\item \textbf{The counting bounds are correct, tight and load-bearing}, and their restatement on
  VCVio's new notation is the same statement (\cref{sec:gen-l0-bounds}).
\item \textbf{The code conforms to the blueprint's Layer 0} except for four points: the explicit
  \lean{n} of \lean{evalOracle_answer}, the unlisted scalar interfaces, an internal inconsistency
  of the blueprint about the $\K$ sampler, and prose that is stale after the upgrade
  (\cref{sec:gen-l0-conformance}).
\end{enumerate}

\subsection{Catalogue}\label{sec:gen-l0-catalogue}

ArkLib's \lean{OracleInterface} (\cref{sec:gen-lib-ark-objects}) is what makes a message an
\emph{oracle}: a query type, and an implementation that answers a query by reading the message.
Its default instance, \lean{instDefault}, answers the query \lean{()} with the whole message.
(\lean{OracleContext ι m} is VCVio's pair of an \lean{OracleSpec ι} and a \lean{QueryImpl} of it
in \lean{m}, \file{VCVio/OracleComp/OracleContext.lean:25-27} at \code{f9dc47d9}.) ArkLib also
registers \lean{instVector : OracleInterface (Vector α n)} answering \emph{positions}
(\file{OracleInterface.lean:333-337}), and \lean{instance (i : Fin 0) : OracleInterface i.elim0}
(\file{:119}).
\begin{leancode}
class OracleInterface (Message : Type u) where
  Query : Type v
  toOC : OracleContext Query (ReaderM Message)
...
def answer {Message : Type*} [O : OracleInterface Message]
    (m : Message) (q : O.Query) : O.Response q :=
  (O.toOC.impl q).run m
...
@[reducible]
def instDefault {Message : Type u} : OracleInterface Message where
  Query := Unit
  toOC.spec := fun _ => Message
  toOC.impl _ := read
\end{leancode}
\src{\file{ArkLib/OracleReduction/OracleInterface.lean:54-57}, \file{:75-78}, \file{:88-96}, ArkLib \code{dca90385}}

{\footnotesize
\begin{longtable}{@{}p{0.30\linewidth}p{0.24\linewidth}p{0.08\linewidth}p{0.27\linewidth}@{}}
\caption{The declarations of Layer 0, at \code{b435631}. Paths are under
\file{LeanerVM/Protocol/}.}\label{tab:gen-l0-catalogue}\\
\toprule
Declaration & File:line & Kind & Plain meaning; what it rests on \\
\midrule
\endfirsthead
\toprule
Declaration & File:line & Kind & Plain meaning; what it rests on \\
\midrule
\endhead
\bottomrule
\endfoot
\lean{Column} & \file{Field.lean:67-69} & structure & a committed column: a table of $2^n$ values
  in $\K$; a \emph{structure} so that ArkLib's position-query \lean{instVector} does not also
  apply. Rests on CompPoly's \lean{CMlPolynomialEval} \\
\lean{finEquivK} & \file{Field.lean:74-78} & def & $\K$ is the $2^{64}$ bit patterns. Core
  \lean{BitVec} \\
\lean{instSampleableTypeK} & \file{Field.lean:81-83} & instance & uniform $\K$: a uniform
  \lean{Fin (2^64)} read as bits. VCVio's \lean{ofEquiv} and \lean{SampleableType (Fin n)} \\
\lean{limbsEquiv} & \file{Field.lean:86-90} & def & $\E$ is its three limbs. CompPoly's
  \lean{Ext} \\
\lean{instSampleableTypeE} & \file{Field.lean:93} & instance & uniform $\E$: three independent
  uniform limbs. VCVio's \lean{Vector} instance, \lean{instSampleableTypeK} \\
\lean{card_E} & \file{Field.lean:96} & theorem & the challenge space has $2^{192}$ elements.
  CompPoly's \lean{card_ext3} \\
\lean{evalOracle} & \file{Field.lean:102-106} & instance & a column as an oracle: query a point
  $r \in \E^n$, get $\mle{q}(r)$. ArkLib's \lean{OracleInterface}, CompPoly's \lean{eval₂Mle} \\
\lean{evalOracle_answer} & \file{Field.lean:109-110} & theorem & the unfolding, for rewriting \\
\lean{instOracleInterfaceE} & \file{Field.lean:115} & instance & a scalar the prover sends: query
  \lean{()}, answer the scalar. ArkLib's \lean{instDefault} \\
\lean{instOracleInterfaceListE} & \file{Field.lean:118} & instance & a list the prover sends:
  query \lean{()}, answer the list. ArkLib's \lean{instDefault} \\
\lean{NoOracle} & \file{ToArkLib/}\allowbreak\file{Oracles.lean:26} & abbrev & the empty family of
  oracles (matches ArkLib's \lean{Fin 0} instance) \\
\lean{OneOracle} & \file{ToArkLib/}\allowbreak\file{Oracles.lean:29} & abbrev & exactly one
  oracle, of type \lean{M} \\
\lean{noOracle_eq} & \file{ToArkLib/}\allowbreak\file{Oracles.lean:32} & theorem & the empty
  family has one inhabitant \\
\lean{probEvent_uniformSample_}\allowbreak\lean{le_of_card_le} &
  \file{ToVCVio/}\allowbreak\file{UniformSample.lean:35-40} & theorem & an event with at most $k$
  witnesses has probability at most $k/\abs{\alpha}$. VCVio's \lean{probEvent_uniformSample} \\
\lean{probEvent_uniformSample_}\allowbreak\lean{le_of_subsingleton} &
  \file{ToVCVio/}\allowbreak\file{UniformSample.lean:44-51} & theorem & an event with at most one
  witness has probability at most $1/\abs{\alpha}$. The previous bound \\
\file{Basic.lean} & \file{Basic.lean} & --- & an empty \code{public section} in
  \lean{namespace LeanerVM.Protocol} (18 lines): nothing \\
\end{longtable}}

The declarations themselves (the counting bounds are in \cref{sec:gen-l0-bounds}):
\begin{leancode}
/-- A column of height `2^n`: the values of a `K`-multilinear on the cube, low bit first. A
structure rather than an abbreviation so that its evaluation oracle below does not overlap
ArkLib's position-query interface on `Vector`, which Layer 11's codeword oracles use. -/
structure Column (n : ℕ) where
  /-- The value table, CompPoly's hypercube representation. -/
  values : CMlPolynomialEval K n

/-! ## Sampling -/

/-- `K` as the `2^64` bit patterns, the `BitVec` view of `Fin (2^64)`. -/
def finEquivK : Fin (2 ^ 64) ≃ K where
  toFun := BitVec.ofFin
  invFun := BitVec.toFin
  left_inv _ := rfl
  right_inv _ := rfl

/-- Uniform sampling of `K`: a uniform `Fin (2^64)` read as a bit pattern. -/
instance instSampleableTypeK : SampleableType K :=
  haveI : NeZero (2 ^ 64) := ⟨by norm_num⟩
  SampleableType.ofEquiv finEquivK

/-- `E` as its three limbs, CompPoly's carrier view. -/
def limbsEquiv : Vector K 3 ≃ E where
  toFun := fun v ↦ CompPoly.Extension.Ext.ofVector (P := BF64.ext3Params) v
  invFun := fun x ↦ CompPoly.Extension.Ext.coeffs (P := BF64.ext3Params) x
  left_inv _ := rfl
  right_inv _ := rfl

/-- Uniform sampling of `E`: three independent uniform limbs. -/
instance instSampleableTypeE : SampleableType E := SampleableType.ofEquiv limbsEquiv

/-- `|E| = 2^192`: the size of the challenge space in every error bound. -/
theorem card_E : Fintype.card E = 2 ^ 192 := BF64.card_ext3

/-! ## The evaluation oracle -/

/-- A column as an oracle: a query is a point `r ∈ E^n`, the answer is the multilinear extension
`q̃(r)`, lifted from `K` to `E`. -/
instance evalOracle (n : ℕ) : OracleInterface (Column n) where
  Query := Vector E n
  toOC :=
    { spec := (Vector E n) →ₒ E
      impl := fun r ↦ do return CMlPolynomialEval.eval₂Mle (← read).values (algebraMap K E) r }

/-- The oracle answers the lifted multilinear extension. -/
theorem evalOracle_answer (n : ℕ) (q : Column n) (r : Vector E n) :
    OracleInterface.answer q r = CMlPolynomialEval.eval₂Mle q.values (algebraMap K E) r := rfl

/-! ## Scalar messages -/

/-- A scalar the prover sends is queried trivially: the answer is the scalar. -/
instance instOracleInterfaceE : OracleInterface E := OracleInterface.instDefault

/-- A list of scalars the prover sends is queried trivially: the answer is the list. -/
instance instOracleInterfaceListE : OracleInterface (List E) := OracleInterface.instDefault
\end{leancode}
\src{\file{LeanerVM/Protocol/Field.lean:64-118}, leanerVM \code{b435631}}

\begin{leancode}
/-- No oracle: the empty family. -/
abbrev NoOracle : Fin 0 → Type := fun i ↦ i.elim0

/-- Exactly one oracle, of type `M`. -/
abbrev OneOracle (M : Type) : Fin 1 → Type := fun _ ↦ M

/-- All functions out of `Fin 0` are equal, so the empty oracle family has one inhabitant. -/
theorem noOracle_eq (o o' : ∀ i, NoOracle i) : o = o' := funext fun i ↦ i.elim0
\end{leancode}
\src{\file{LeanerVM/Protocol/ToArkLib/Oracles.lean:25-32}, leanerVM \code{b435631}}

\paragraph{Axioms.} Every declaration above depends on at most \lean{propext},
\lean{Classical.choice}, \lean{Quot.sound} (\lean{noOracle_eq}: \lean{Quot.sound} only;
\lean{finEquivK}: \lean{propext} only; probe \code{Layer0}, section 3).

\paragraph{Tests.} \file{tests/LeanerVMTests/Protocol/Field.lean} (at \code{b435631}): a
two-variable column \texttt{⟨\#v[1, 2, 3, 4]⟩} (\file{:27}); four \texttt{\#guard}s that the oracle
returns the table entry on the cube (\file{:33-36}); one off the cube,
\texttt{answer col \#v[y, ofK 0] = ofK 1 + ofK 3 * y} (\file{:39}, where at this pin $2$ and $3$ are
bit patterns, so $1 + y\cdot(1 \oplus 2)$); agreement with \lean{evalMle} of the lifted table at
$(y, y^2)$ (\file{:41-42}); a mutated column answers differently (\file{:44}); \lean{sampleK},
\lean{sampleE} as compiled definitions (\file{:49}, \file{:52}: they fail to compile if a sampler
has no compiler IR); \texttt{\#synth OracleInterface E} and \lean{(List E)} (\file{:59-60});
\lean{example : Fintype.card E = 2 ^ 192 := card_E} (\file{:64}).

\subsection{The samplers: uniform, proved, and the diamond on \texorpdfstring{$\K$}{K}}\label{sec:gen-l0-samplers}

\paragraph{Uniformity is proved, by construction.} A \lean{SampleableType} instance cannot be
non-uniform (\cref{sec:gen-lib-vcvio-sample}: the class carries \lean{probOutput_selectElem_eq}
and \lean{mem_support_selectElem}, and VCVio derives
$\Pr[= x \mid \$^t \alpha] = \abs{\alpha}^{-1}$ for every instance). Both instances are
\lean{SampleableType.ofEquiv} of a Mathlib \lean{Equiv} (\lean{α ≃ β} bundles both inverse laws,
so it is a bijection; here both laws are \lean{rfl}), and \lean{ofEquiv} re-proves the two laws
(\file{SampleableType.lean:260-267} at \code{f9dc47d9}). The chain for $\E$ is: the primitive
coin \lean{$[0..2^64−1]} (\lean{unifSpec.query}, uniform by \lean{IsUniformSpec unifSpec}) $\to$
\lean{finEquivK} $\to$ three independent draws (VCVio's \lean{Vector} instance) $\to$
\lean{limbsEquiv}. The probe \code{Layer0} (section 4) has all of the following accepted:
\begin{leancode}
example : Function.Bijective finEquivK := finEquivK.bijective
example : Function.Bijective limbsEquiv := limbsEquiv.bijective
example (x : K) : Pr[= x | $ᵗ K] = (Fintype.card K : ℝ≥0∞)⁻¹ := probOutput_uniformSample K x
example (x : E) : Pr[= x | $ᵗ E] = (Fintype.card E : ℝ≥0∞)⁻¹ := probOutput_uniformSample E x
example (x : E) : Pr[= x | $ᵗ E] = ((2 ^ 192 : ℕ) : ℝ≥0∞)⁻¹ := by
  rw [probOutput_uniformSample, card_E]
example : Pr[⊥ | $ᵗ E] = 0 := probFailure_uniformSample E
example : support ($ᵗ E) = Set.univ := support_uniformSample E
example : ($ᵗ E) = limbsEquiv <$> ($ᵗ (Vector K 3)) := rfl
example : ($ᵗ K) = finEquivK <$> ($ᵗ (Fin (2 ^ 64))) := rfl
\end{leancode}
\src{probe \code{Layer0} (\file{probes/lib-others/Layer0.lean}), section 4, run at the old pins}

Instance search (probe \code{Layer0}, section 1): \lean{SampleableType K} $\to$
\lean{instSampleableTypeK}, \lean{SampleableType E} $\to$ \lean{instSampleableTypeE},
\lean{SampleableType (Vector K 3)} $\to$ \lean{instSampleableTypeVector K 3},
\lean{SampleableType (Fin (2 ^ 64))} $\to$ \lean{instSampleableTypeFinOfNeZeroNat (2 ^ 64)}.

\paragraph{The diamond on \texorpdfstring{$\K$}{K} (old pins).} Because $\K$ is \lean{BitVec 64}
to instance search and Mathlib has \lean{FinEnum (BitVec n)}, VCVio's
\lean{FinEnum.SampleableType} already supplies a sampler on $\K$ \emph{without} leanerVM's
instance. The probe \code{SamplerDiamond} (which imports \lean{LeanerVM.Parameters.Field} and
VCVio only) shows it:
\begin{leancode}
#synth SampleableType K        ⟹  FinEnum.SampleableType K
#synth SampleableType (BitVec 64)  ⟹  FinEnum.SampleableType (BitVec 64)
#synth FinEnum K               ⟹  FinEnum.instBitVec 64
#synth SampleableType E        ⟹  error: failed to synthesize SampleableType E
\end{leancode}
\src{probe \code{SamplerDiamond} (\file{probes/lib-others/SamplerDiamond.lean}), run at the old pins; the last line is an intended error}

With \lean{LeanerVM.Protocol.Field} imported, leanerVM's instance is found (declared later, same
priority), and the two are \textbf{the same term up to unfolding}; the probe
\code{SamplerDiamond2} has both of these accepted (exit 0):
\begin{leancode}
example :
    (instSampleableTypeK).selectElem = (FinEnum.SampleableType K).selectElem := rfl
example : (instSampleableTypeK : SampleableType K) = FinEnum.SampleableType K := rfl
\end{leancode}
\src{probe \code{SamplerDiamond2} (\file{probes/lib-others/SamplerDiamond2.lean}), run at the old pins}

In any case any two samplers on a finite type give every element the same probability (probe
\code{SamplerDiamond}, accepted on the standard axioms):
\begin{leancode}
theorem sampler_unique {α : Type} [Fintype α] (s t : SampleableType α) (x : α) :
    Pr[= x | @uniformSample α s] = Pr[= x | @uniformSample α t] := by
  rw [@probOutput_uniformSample α s, @probOutput_uniformSample α t]
\end{leancode}
\src{probe \code{SamplerDiamond} (\file{probes/lib-others/SamplerDiamond.lean}), run at the old pins}

So the diamond is harmless: no statement can depend on which instance is chosen. On $\E$ there is
exactly one instance in scope. \textbf{Both samplers are uniform, proved; no distribution
question remains.}

\paragraph{\texorpdfstring{\lean{SampleableType K}}{SampleableType K} has no consumer.} No
declaration outside \file{Protocol/Field.lean} mentions it (a \code{git grep} of \code{b435631}
over \file{LeanerVM}, \file{tests} and \file{docs} finds only the blueprint's Layer 0 sketch), and
leanVM samples every challenge in $\E$ (\file{crates/fiat_shamir/src/lib.rs:95-99},
\code{pub fn sample(&mut self) -> F192}; the module docstring
\file{crates/primitives/src/field/mod.rs:3-6}: \enquote{challenges, sumcheck/GKR values, and
transcript scalars are E-valued}). The blueprint's own dependency row says \enquote{Layer 0
supplies the \lean{E} instance} (\file{protocol-blueprint.md:251}). At the old pins the instance
is moreover a definitional duplicate of the library one. It is still needed as a \emph{building
block} of \lean{instSampleableTypeE}, through VCVio's \lean{Vector} instance; at the old pin the
library \lean{FinEnum} instance would serve. Finding: \cref{f:lib-k-sampler-unused}.

\paragraph{At \texorpdfstring{\code{144c5aa}}{144c5aa} (new pins).} $\K$ is a structure
(\cref{sec:gen-lib-comppoly-new}), so \lean{FinEnum (BitVec 64)} no longer applies to it and the
diamond disappears; \lean{finEquivK} becomes
\lean{toFun := fun i ↦ BF64.ofBitVec (BitVec.ofFin i)}, \lean{invFun := fun x ↦ x.toBitVec.toFin}
(\code{git diff b435631 144c5aa -- LeanerVM/Protocol/Field.lean}); \lean{instSampleableTypeK} is
then the only $\K$ sampler (by reading; not run). VCVio \code{a4232d08} states uniformity directly
as a class law and proves sampler irrelevance upstream, the statement of the probe's
\lean{sampler_unique} (\cref{sec:gen-lib-vcvio-new}).

\subsection{The oracle interface of the stack}\label{sec:gen-l0-oracle}

\begin{itemize}
\item \textbf{Query}: \lean{Vector E n}, a point $r \in \E^n$. \textbf{Answer}: $\E$
  (\lean{spec := (Vector E n) →ₒ E}). \textbf{Implementation}: read the column and return
  \lean{eval₂Mle q.values (algebraMap K E) r}, which is
  \lean{evalMle (map (algebraMap K E) q) r}
  $= \sum_i \mathrm{ofK}(q_i) \cdot \prod_j (\mathrm{bit}_j(i)\ ?\ r_j : 1 - r_j)$, the
  multilinear extension $\mle{q}(r)$ of the $\K$-table lifted to $\E$
  (\cref{sec:gen-lib-comppoly-mle}). The probe \code{Layer0}, section 5, has these accepted:
\begin{leancode}
example (n : ℕ) : OracleInterface.Query (Column n) = Vector E n := rfl
example (n : ℕ) (r : Vector E n) : OracleInterface.Response (Message := Column n) r = E := rfl
example (n : ℕ) (q : Column n) (r : Vector E n) :
    OracleInterface.answer q r =
      CMlPolynomialEval.eval (CMlPolynomialEval.map (algebraMap K E) q.values) r := by
  rw [evalOracle_answer, CMlPolynomialEval.eval₂_mle_eq_eval₂]; rfl
\end{leancode}
\src{probe \code{Layer0} (\file{probes/lib-others/Layer0.lean}), section 5, run at the old pins}
\item This is leanVM's oracle: \enquote{we commit to a $\K$-valued multilinear, and opening claims
  are in $\E$} (\file{doc/leanvm/body/b-polynomial-commitment-scheme.tex:4}).
\item \textbf{Name}: the code calls it \lean{evalOracle}, as the blueprint does
  (\file{protocol-blueprint.md:638}); its full name is \lean{LeanerVM.Protocol.evalOracle} (probe
  \code{Layer0}: \texttt{\#synth OracleInterface (Column 3)} $\Rightarrow$ \lean{evalOracle 3}).
\item \textbf{Why a structure.} \lean{CMlPolynomialEval K n} is \lean{@[reducible]}
  \lean{Vector K (2^n)}, and ArkLib has \lean{instVector : OracleInterface (Vector α n)} with
  \lean{Query := Fin n} (positions). Were \lean{Column} an abbreviation, a column could be queried
  at positions. The structure prevents it; the justification is correct at \code{dca90385}
  (\file{OracleInterface.lean:333-337}).
\item \textbf{Other instances.} \lean{OracleInterface E} and \lean{OracleInterface (List E)} are
  ArkLib's \lean{instDefault}: query \lean{Unit}, answer the whole message (probe \code{Layer0}:
  \lean{OracleInterface.Query E = Unit := rfl}, \lean{OracleInterface.answer x () = x := rfl}, the
  same for \lean{List E}). They are the \enquote{the verifier reads what the prover sent} interface
  for scalars and coefficient lists; ArkLib registers no default instance for any type, at
  \code{dca90385} or at \code{7653a901} (a \code{git grep} for \lean{instDefault} instances finds
  only two local ones in \file{ArkLib/ProofSystem/RingSwitching/Packing/Spec.lean:105-106} at
  \code{7653a901}), so there is no ambiguity yet. Finding: \cref{f:lib-scalar-interfaces}.
\end{itemize}
No issue found with the oracle interface: the query space, the answer and the lift are those of
the specification, and the mutated-column test makes the answer depend on the column.

\subsection{\texorpdfstring{\lean{card_E}}{card\_E}}\label{sec:gen-l0-card}

\lean{theorem card_E : Fintype.card E = 2 ^ 192 := BF64.card_ext3} (\file{Protocol/Field.lean:96}).
The \lean{Fintype E} instance is CompPoly's \lean{Ext.instFintype} (probe \code{Layer0}:
\texttt{\#synth Fintype E} $\Rightarrow$ \lean{Extension.Ext.instFintype}), and \lean{card_ext3} is
$(2^{64})^3$ by the coefficient bijection (\cref{sec:gen-lib-comppoly-fields}). \lean{Fintype.card}
does not depend on the instance (all \lean{Fintype α} instances are equal: \lean{Fintype} is a
subsingleton), so the statement is robust. Correct. At \code{144c5aa} the proof is unchanged in
text; \lean{Fintype K} is leanerVM's proof-only \lean{Fintype.ofFinite K} and
\lean{card_ext3 [Fintype Ext3]} accepts any instance (\cref{sec:gen-lib-comppoly-new}).

\subsection{The counting bounds}\label{sec:gen-l0-bounds}

\begin{leancode}
variable {α : Type} [SampleableType α] [Fintype α]

/-- An event with at most `k` witnesses has probability at most `k/|α|` under a uniform
sample. -/
theorem probEvent_uniformSample_le_of_card_le (p : α → Prop) [DecidablePred p] {k : ℕ}
    (h : (Finset.univ.filter p).card ≤ k) :
    Pr[p | $ᵗ α] ≤ ((k / Fintype.card α : ℝ≥0) : ℝ≥0∞) := by
...
/-- An event any two of whose witnesses are equal has probability at most `1/|α|` under a
uniform sample. -/
theorem probEvent_uniformSample_le_of_subsingleton (p : α → Prop)
    (h : ∀ a b, p a → p b → a = b) :
    Pr[p | $ᵗ α] ≤ ((1 / Fintype.card α : ℝ≥0) : ℝ≥0∞) := by
\end{leancode}
\src{\file{LeanerVM/Protocol/ToVCVio/UniformSample.lean:31-51} (proofs elided), leanerVM \code{b435631}}

\textbf{Correct.} The left side is $\abs{\mathrm{filter}\ p} / \abs{\alpha}$ exactly (VCVio
\lean{probEvent_uniformSample}, \cref{sec:gen-lib-vcvio-sample}); the right side is the real
$k/\abs{\alpha}$ (division in \lean{ℝ≥0}, then coerced; $\abs{\alpha} \neq 0$ because a
\lean{SampleableType} is nonempty). \textbf{Not vacuous, and the hypotheses carry the weight.}
The probe \code{CountingBounds} has all of the following accepted, on the standard axioms (the
only messages are unused-variable warnings): the single-witness bound is attained, the
two-witness event exceeds it, and the bound with $k = 2$ is attained.
\begin{leancode}
theorem one_witness : Pr[fun x : Fin 4 ↦ x = 2 | $ᵗ (Fin 4)] = 1 / 4
theorem one_witness_bound :
    Pr[fun x : Fin 4 ↦ x = 2 | $ᵗ (Fin 4)] ≤ ((1 / Fintype.card (Fin 4) : ℝ≥0) : ℝ≥0∞)
theorem rhs_eq : ((1 / Fintype.card (Fin 4) : ℝ≥0) : ℝ≥0∞) = 1 / 4          -- attained
theorem two_witnesses : Pr[fun x : Fin 4 ↦ x = 1 ∨ x = 2 | $ᵗ (Fin 4)] = 1 / 2
theorem two_witnesses_exceed :
    ¬ Pr[fun x : Fin 4 ↦ x = 1 ∨ x = 2 | $ᵗ (Fin 4)] ≤ 1 / 4               -- hypothesis needed
theorem two_witnesses_bound :
    Pr[fun x : Fin 4 ↦ x = 1 ∨ x = 2 | $ᵗ (Fin 4)]
      ≤ (((2 : ℕ) / Fintype.card (Fin 4) : ℝ≥0) : ℝ≥0∞)                     -- k = 2, attained
theorem one_bad_challenge (c : E) :
    Pr[fun x : E ↦ x = c | $ᵗ E] = ((2 ^ 192 : ℕ) : ℝ≥0∞)⁻¹
theorem one_bad_challenge_bound (c : E) :
    Pr[fun x : E ↦ x = c | $ᵗ E] ≤ ((1 / Fintype.card E : ℝ≥0) : ℝ≥0∞)
\end{leancode}
\src{probe \code{CountingBounds} (\file{probes/lib-others/CountingBounds.lean}), run at the old pins}

Consumers at \code{b435631}: \lean{…_of_subsingleton} in the public-input phase
(\file{PublicInput.lean:362}); \lean{…_of_card_le} only through it.

\paragraph{At \texorpdfstring{\code{144c5aa}}{144c5aa}.} Both statements were restated on VCVio's
new probability API (\code{git diff b435631 144c5aa -- LeanerVM/Protocol/ToVCVio/UniformSample.lean}):
only the left side changes, \lean{Pr[p | $ᵗ α]} becomes \lean{Pr{let sample ← $ᵗ α}[p sample]};
binders, hypotheses and right sides are unchanged. At VCVio \code{a4232d08}, \lean{Pr{…}[…]} is a
macro for the measure of the event \lean{{True}} under the computation returning the proposition
(\cref{sec:gen-lib-vcvio-new}). \textbf{The two statements are the same mathematical statement.}
At each pin the left side equals the same closed form by a library theorem: at \code{f9dc47d9},
\lean{probEvent_uniformSample : Pr[ p | $ᵗ α] = (Finset.univ.filter p).card / Fintype.card α}
(\file{SampleableType.lean:224-226}); at \code{a4232d08},
\lean{SampleableType.prEvent_uniformSample : Pr{let x ← $ᵗ α}[p x] = (Finset.univ.filter p).card / (Fintype.card α : ENNReal)}
(\file{Constructions/SampleableType/NativeMeasure.lean:41-45}). And VCVio \code{a4232d08} proves
its measure semantics agrees with the old mass-function semantics on discrete answer spaces. So
the restated bounds bound the same number by the same number (verified by reading; not re-run at
the new pin).

\paragraph{Upstream at the new pin.} At VCVio \code{f9dc47d9} there is no duplicate: VCVio has
the equalities \lean{probEvent_uniformSample} and \lean{probOutput_uniformSample}, and the two
local bounds are their two-line corollaries. At VCVio \code{a4232d08},
\lean{…_of_card_le} is, up to a coercion, the $\Leftarrow$ direction of the following theorem
(section variables \lean{{α : Type} [SampleableType α] [Fintype α] (p : α → Prop) [DecidablePred p]},
\file{:117}); the local right side \lean{((k / Fintype.card α : ℝ≥0) : ℝ≥0∞)} equals
\lean{k / (Fintype.card α : ℝ≥0∞)} by \lean{ENNReal.coe_div}.
\begin{leancode}
theorem prEvent_uniformSample_le_div_iff {c : ℕ} :
    Pr{let x ← $ᵗ α}[p x] ≤ c / (Fintype.card α : ℝ≥0∞) ↔
      (Finset.univ.filter p).card ≤ c
\end{leancode}
\src{\file{VCVio/OracleComp/Constructions/SampleableType/NativeMeasure.lean:126-130}, VCVio \code{a4232d08}}

For single-witness events VCVio has
\lean{prEvent_uniformSample_eq_singleton : Pr{let x ← $ᵗ α}[x = a] = (Fintype.card α : ℝ≥0∞)⁻¹}
(\file{:110-111}), but no \enquote{at most one witness} form (a search for \code{subsingleton} and
\code{card_le_one} in \file{SampleableType/} and \file{EvalDist/ProbabilityBounds.lean} finds
none). The probe's \lean{sampler_unique} is upstream as \lean{prEvent_uniformSample_inst_irrel}
(\file{:280-284}). So at the new pin \lean{…_of_card_le} is redundant and can be replaced by the
upstream equivalence; \lean{…_of_subsingleton} is kept, or upstreamed
(\cref{f:lib-counting-bound-redundant}).

\subsection{Conformance with the blueprint's Layer 0}\label{sec:gen-l0-conformance}

The blueprint's Layer 0 is \file{protocol-blueprint.md:625-650} at \code{b435631}.

{\footnotesize
\begin{longtable}{@{}p{0.38\linewidth}p{0.28\linewidth}p{0.26\linewidth}@{}}
\caption{The blueprint's Layer 0 against the code at \code{b435631}.}\label{tab:gen-l0-conformance}\\
\toprule
Blueprint & Code & Verdict \\
\midrule
\endfirsthead
\toprule
Blueprint & Code & Verdict \\
\midrule
\endhead
\bottomrule
\endfoot
\code{instance : SampleableType K  -- a uniform Fin (2^64) read as a bit pattern}
  & \lean{instSampleableTypeK := SampleableType.ofEquiv finEquivK}
  & agrees; unused (\cref{sec:gen-l0-samplers}) \\
\code{instance : SampleableType E  -- three independent limbs through Ext.ofVector; never enumerates}
  & \lean{instSampleableTypeE := SampleableType.ofEquiv limbsEquiv}
  & agrees \\
\code{theorem card_E : Fintype.card E = 2 ^ 192}
  & identical & agrees \\
\code{instance evalOracle (n : ℕ) : OracleInterface (Column n) where Query := Vector E n … answer := eval₂Mle q (algebraMap K E)}
  & identical & agrees \\
\code{theorem evalOracle_answer (q : Column n) (r : Vector E n) : …}
  & \code{theorem evalOracle_answer (n : ℕ) (q : Column n) (r : Vector E n) : …}
  & \textbf{mismatch}: \lean{n} explicit in the code, implicit (auto-bound) in the sketch;
    callers must write \lean{evalOracle_answer n q r} although \lean{n} is determined by \lean{q}
    (\cref{f:lib-evaloracle-answer-n}) \\
\enquote{\lean{SampleableType.ofEquiv} on the carrier views, as ArkLib's own
  \lean{KoalaBear.Ext6.sampleableType}, and a test probes that it has compiler IR}
  & yes (\lean{sampleK}, \lean{sampleE}) & agrees \\
Tests: \enquote{\texttt{\#guard}s that the oracle answers \lean{evalMle} on a two-variable table, on
  and off the cube, with a mutated column; \lean{example : Fintype.card E = 2 ^ 192 := card_E}}
  & yes (\cref{sec:gen-l0-catalogue}) & agrees \\
---
  & \lean{instOracleInterfaceE}, \lean{instOracleInterfaceListE} in \file{Protocol/Field.lean}
  & \textbf{not in the Layer 0 sketch}, nor in the interface list \file{protocol-blueprint.md:1400}
    (\enquote{\code{instSampleableTypeE  evalOracle  card_E}}), which also omits
    \lean{instSampleableTypeK} and \lean{Column} (the latter is at \file{:666})
    (\cref{f:lib-scalar-interfaces}) \\
\file{protocol-blueprint.md:251}: \enquote{Layer 0 supplies the \lean{E} instance}
  & Layer 0 supplies $\K$ and $\E$
  & internal inconsistency of the blueprint (\cref{f:lib-k-sampler-unused}) \\
\enquote{Add ArkLib at \code{dca90385} \dots\ the root pin \code{3468b38c} wins}
  (\file{:629-631}; \file{:633-635} at \code{144c5aa}), \enquote{The sampler is the point at which
  a wrong \lean{Fintype} would enumerate $2^{64}$ or $2^{192}$ elements (finding P3)}
  (\file{:645-646}; \file{:649-650} at \code{144c5aa})
  & at \code{144c5aa} the pins moved and \lean{Fintype K} is proof-only
  & stale after \#61 (the disclaimer added at \file{144c5aa:protocol-blueprint.md:52-54} covers
    \enquote{dependency API tables and signatures}, not this prose)
    (\cref{f:lib-layer0-stale}) \\
\end{longtable}}

\lean{NoOracle} and \lean{OneOracle} are specified in the spine section
(\file{protocol-blueprint.md:472}, \file{:1377}), and \file{ToVCVio/UniformSample.lean}'s bound is
cited at \file{:1046}; both agree with the code.

\par\endgroup